\chapter{Recurrence CFD: Simulasi dari Ingatan}
\label{ch:rcfd}

Recurrence CFD adalah topik tugas akhir saya, dan kuliah metode numerik fluida memberinya satu pertemuan
penuh. Stefan Pirker membukanya dengan Max dan Moritz, dua anak nakal dari cerita bergambar Wilhelm
Busch, dan guru mereka, Lehrer Lämpel. Max dan Moritz membuat tebakan yang berani, sang guru
mengoreksinya. Menurut Pirker, begitulah cara algoritme SIMPLE menyelesaikan persamaan Navier--Stokes:
tebakan yang tidak sempurna, lalu koreksi berbasis fisika. Dan begitulah juga rCFD dikembangkan.
Bab ini mengikuti alur kuliah itu, lalu menceritakan sedikit tentang apa yang saya kerjakan dalam tugas
akhir.

\section{Aliran yang mengulang dirinya}

Banyak aliran yang menarik bagi industri berjalan lama sekali. Peleburan baja di dalam ladle,
pencampuran di tangki berpengaduk, dan gelembung yang naik di kolom gelembung bisa berlangsung menit
sampai jam, sedangkan simulasi CFD yang teliti hanya sanggup menjangkau beberapa puluh detik waktu fisis
dalam waktu hitung yang wajar. Untungnya, banyak aliran seperti ini, setelah melewati fase awal, tidak
menghasilkan pola yang benar-benar baru. Pusaran datang dan pergi, gelembung naik dalam pola yang
mirip, dan keadaan aliran berulang kali kembali ke dekat keadaan yang pernah dikunjungi sebelumnya.

Pengulangan ini bisa diukur. Simpan medan aliran dari simulasi penuh sebagai deretan bingkai
$\bm\phi_1,\dots,\bm\phi_M$. Untuk setiap pasangan bingkai, hitung jaraknya, lalu normalisasikan:
\begin{equation}
R_{ij}=\frac{\|\bm\phi_i-\bm\phi_j\|}{\max_{k,l}\|\bm\phi_k-\bm\phi_l\|}.
\end{equation}
Matriks $\bm R$ ini disebut \istilah{recurrence matrix}. Diagonalnya nol, karena setiap bingkai identik
dengan dirinya sendiri. Kalau aliran periodik sempurna, garis-garis gelap sejajar diagonal muncul pada
jarak satu periode. Pada aliran turbulen, polanya lebih kabur, tetapi tetap ada daerah-daerah di luar
diagonal dengan nilai kecil: pasangan bingkai yang berjauhan dalam waktu tetapi mirip satu sama lain.
Dalam bahasa machine learning, $\bm R$ adalah matriks jarak antar sampel, dan pasangan dengan jarak kecil
adalah tetangga terdekat di ruang keadaan aliran. \Cref{fig:recurrence} memperlihatkan bentuk khas matriks ini
untuk aliran yang hampir periodik.

\begin{figure}[ht]
\centering
\begin{tikzpicture}
\begin{axis}[width=0.5\linewidth,height=0.5\linewidth,view={0}{90},axis on top,xmin=0,xmax=60,ymin=0,ymax=60,
  xlabel={\small bingkai $i$},ylabel={\small bingkai $j$},xtick={0,20,40,60},ytick={0,20,40,60},
  colormap={abu}{color=(biru!85!black) color=(biru!25) color=(pasir!40)},y dir=reverse,
  colorbar,colorbar style={width=2.5mm,ytick={0,1},ylabel={\small $R_{ij}$}}]
  \addplot3[surf,shader=interp,samples=61,domain=0:60,y domain=0:60]
    {min(1,abs(sin(deg(pi*(x-y)/20)))*(1+0.004*abs(x-y)))};
\end{axis}
\end{tikzpicture}
\caption{Bentuk recurrence matrix untuk aliran yang hampir periodik dengan periode sekitar dua puluh bingkai (ilustrasi).
Garis gelap di diagonal dan sejajar diagonal menandai pasangan bingkai yang mirip, kandidat tempat melompat saat
menjahit aliran.}
\label{fig:recurrence}
\end{figure}

\section{rCFD berbasis aliran}

Gagasan pertama, yang diterbitkan \citet{lichtenegger2016} untuk kolom gelembung, berbunyi seperti
kalimat yang berulang di slide: ``kita putuskan dengan berani untuk menyusun aliran buatan dengan
menjahit potongan-potongan rangkaian aliran''. Mulailah dari sebuah bingkai, mainkan beberapa bingkai
berikutnya dari basis data, lalu lompat ke bingkai lain yang mirip dengan bingkai terakhir menurut
recurrence matrix, dan teruskan dari sana. Hasilnya adalah deret medan aliran sepanjang yang kita mau,
yang secara statistik meniru aliran asli, tanpa pernah menyelesaikan persamaan Navier--Stokes lagi.

Di atas aliran yang dijahit ini, besaran yang terbawa aliran, misalnya konsentrasi zat terlarut, dilacak
dengan persamaan transport implisit (\cref{eq:transport}). Persamaan transport untuk satu besaran skalar
jauh lebih murah daripada persamaan Navier--Stokes yang terkopel, dan karena implisit, langkah waktunya
bisa jauh lebih besar. Pada kasus kolom gelembung, rCFD berbasis aliran menghasilkan prediksi yang cocok
dengan CFD penuh dengan percepatan sekitar seratus kali.

Bayangkan mencoba meramalkan arus lalu lintas sebuah kota selama setahun. Daripada mensimulasikan setiap
mobil, kita merekam lalu lintas selama beberapa minggu, lalu menyusun tahun buatan dengan menyambung
potongan-potongan hari yang mirip: setelah Senin pagi yang macet biasanya datang siang yang agak sepi.
Kalau yang kita pedulikan adalah ke mana sebuah paket akan sampai, rekaman lalu lintas yang dijahit sudah
cukup.

\section{rCFD berbasis transport}

Langkah kedua, dari \citet{pirker2018}, bahkan melewati persamaan transport. Kuliah menjelaskannya dengan
sungai Danube di Linz. Sebuah zat yang dijatuhkan ke sungai dibawa arus dan menyebar. Arus membawa zat
dari satu tempat ke tempat lain, yaitu konveksi, dan pusaran kecil membuatnya menyebar, yaitu difusi.
Dalam bahasa sel-sel volume hingga, konveksi berarti sebagian isi sebuah sel berpindah ke sel lain, dan
difusi berarti dua sel bertetangga saling bertukar isi. Mesh simulasi adalah jaringan komunikasi antar
sel. rCFD berbasis transport memutuskan, sekali lagi dengan berani, untuk mengganti konveksi dengan
\istilah{cell-to-cell shifts} dan difusi dengan \istilah{face swaps} yang direkam dari simulasi penuh.

Secara sederhana, satu langkah transport menjadi perkalian dengan sebuah matriks pemindahan. Kalau $c_i$
adalah jumlah zat di sel $i$ dan $P_{ji}$ fraksi isi sel $i$ yang dipindahkan ke sel $j$ dalam satu
langkah, maka
\begin{equation}
c_j^{\mathrm{baru}}=\sum_iP_{ji}\,c_i.
\end{equation}
Massa total kekal tepat ketika setiap kolom $\bm P$ berjumlah satu, $\sum_jP_{ji}=1$: semua isi sel $i$
pergi ke suatu tempat, tidak ada yang hilang dan tidak ada yang muncul. Buktinya satu baris: jumlah semua $c_j^{\mathrm{baru}}$
adalah $\sum_i c_i\sum_jP_{ji}=\sum_ic_i$.

Contoh tiga sel di sebuah saluran tertutup membuatnya konkret. Aliran membawa isi dari sel pertama ke sel kedua dan
dari sel kedua ke sel ketiga, dan sel ketiga adalah ujung buntu:
\begin{equation}
\bm P=\begin{pmatrix}0{,}7&0&0\\0{,}3&0{,}6&0\\0&0{,}4&1\end{pmatrix}.
\end{equation}
Setiap kolom berjumlah satu. Mulai dengan seluruh zat di sel pertama, $\bm c=(1,0,0)$. Setelah satu langkah
$\bm c=(0{,}7;\ 0{,}3;\ 0)$, setelah dua langkah $(0{,}49;\ 0{,}39;\ 0{,}12)$, dan jumlahnya tetap satu. Sekarang bayangkan
kolom pertama hanya berjumlah $0{,}98$ karena dua bingkai yang dijahit tidak cocok sempurna. Zat yang melewati sel itu
kehilangan dua persen setiap kali, dan setelah seratus langkah, $0{,}98^{100}\approx0{,}13$. Galat yang tampak kecil
dalam satu langkah menghabiskan hampir seluruh zat dalam simulasi panjang. Inilah bentuk kegagalan pertama di bawah,
dan juga alasan mengapa kekekalan harus dijaga secara struktural, bukan diharapkan muncul sendiri. Pembaca yang ingat \cref{ch:penemuan}
akan mengenali bahwa matriks seperti ini adalah versi diskret dari operator transfer, saudara operator
Koopman, yang memajukan distribusi alih-alih memajukan pengamatan.

Selisih kecepatannya besar sekali. Untuk simulasi pencampuran di sebuah ladle baja, slide kuliah membandingkan CFD
penuh yang membutuhkan 26 jam dengan rCFD berbasis transport yang selesai dalam 1,6 detik pada perangkat
keras yang sama, dengan histogram konsentrasi yang cocok di beberapa titik waktu.

\section{Memilih lompatan dan menilai hasil}

Penjahitan aliran membutuhkan dua keputusan praktis. Pertama, kapan melompat. Kalau lompatan dilakukan terlalu sering,
aliran yang dijahit kehilangan kontinuitas waktu, dan struktur yang seharusnya terbawa beberapa langkah terputus di
tengah jalan. Kalau terlalu jarang, basis data dimainkan ulang hampir apa adanya, dan statistik jangka panjangnya hanya
sebaik panjang rekaman. Biasanya rekaman dimainkan selama beberapa bingkai berturut-turut, lalu lompatan dilakukan ke salah
satu dari beberapa kandidat terdekat menurut recurrence matrix, dipilih secara acak agar urutan yang dijahit tidak berputar
di lingkaran yang sama.

Kedua, berapa jauh boleh melompat. Ambang kemiripan yang ketat menjamin sambungan yang mulus tetapi mempersempit pilihan,
sedangkan ambang longgar memberi banyak pilihan dengan sambungan yang lebih kasar. Masalah ini mirip dengan pemilihan $k$
pada $k$-nearest neighbour di \cref{ch:unsup}, dan jawabannya juga mirip: tidak ada angka yang benar untuk semua kasus,
sehingga pengaruhnya harus diuji.

Menilai hasil rCFD juga butuh hati-hati. Karena aliran yang dijahit tidak pernah identik dengan aliran asli bingkai demi
bingkai, membandingkan medan sesaat tidak adil, alasan yang sama dengan evaluasi rollout di \cref{ch:operator}. Yang
dibandingkan adalah besaran yang dipedulikan: kurva konsentrasi atau temperatur di titik-titik pemantauan, waktu yang
dibutuhkan untuk mencapai keadaan tercampur, dan distribusi besaran di seluruh domain. Neraca global, seperti total massa
atau energi, menjadi pemeriksaan pertama yang murah, karena galat neraca yang tumbuh adalah tanda paling awal bahwa
sesuatu di dalam matriks pemindahan tidak beres.

\section{Tiga kegagalan berani dan koreksinya}

Bagian paling jujur dari kuliah adalah tiga ``kegagalan berani'' rCFD dan cara memperbaikinya.

Kegagalan pertama adalah ketidakseimbangan massa. Untuk simulasi yang panjang, hasil rCFD menyimpang
dari CFD penuh. Penyebabnya, pemindahan antar sel yang direkam dari bingkai berbeda dan dijahit bersama
tidak menjamin bahwa setiap kolom $\bm P$ berjumlah tepat satu. Secara lokal, sebagian massa hilang dan
sebagian lain tercipta. Secara global, neraca massa spesies bisa dihitung dengan mudah, sehingga galat
rCFD bisa dipantau. Koreksinya memakai pemantau global itu untuk menyesuaikan pola konsentrasi lokal
lewat face swaps yang bekerja pada massa, sehingga neraca kembali seimbang.

Kegagalan kedua adalah difusi yang tidak fisis. Face swaps yang sama kuatnya di semua tempat membuat zat
menyebar terlalu cepat di daerah tenang dan terlalu lambat di daerah turbulen. \citet{du2020}, yang
menerapkan rCFD pada penyebaran polutan di sekitar bangunan, memperbaikinya dengan membuat kekuatan face
swaps bergantung pada turbulensi lokal, yang dalam LES dinyatakan oleh viskositas sub-grid.

Kegagalan ketiga adalah pengulangan yang kurang menonjol. Di domain yang besar, tidak ada dua bingkai
yang mirip di \emph{semua} tempat sekaligus. Penjahitan global memasangkan bingkai yang secara rata-rata
mirip tetapi secara lokal, misalnya di ekor pusaran di belakang sebuah gedung, tidak mirip sama sekali.
Usulan koreksinya adalah memecah domain menjadi ``pulau-pulau'' dengan pengulangan yang kuat dan ``laut''
latar dengan pengulangan yang lemah, masing-masing dengan recurrence matrix sendiri.

\section{Menuju kembaran digital}

Kuliah ditutup dengan visi yang lebih besar. Kalau rCFD begitu cepat, ia bisa berjalan bersamaan dengan
proses industri yang sebenarnya, sebagai kembaran digital. Masalahnya, proses nyata punya lebih dari satu
keadaan operasi, dan rCFD yang dijalankan maju tanpa koreksi pasti melenceng dari kenyataan. Jalan
keluarnya mirip asimilasi data di \cref{ch:prob}: jalankan rCFD sambil mengoreksinya dengan pengukuran
proses yang tersedia secara langsung. Kuliah memperlihatkan percobaan sederhana untuk ide ini, dua aliran
air panas dan dingin yang dicampur di sebuah tangki dengan deretan sensor temperatur vertikal. Model
tangki ideal yang tercampur sempurna terlalu sederhana karena mengabaikan perbedaan di dalam tangki,
sedangkan CFD tanpa validasi rawan galat. rCFD yang diturunkan dari CFD dan dikoreksi dengan data sensor
menempati tempat di antara keduanya.

\section{Tugas akhir saya}

Tugas akhir saya menerapkan rCFD berbasis transport pada pancaran apung (\istilah{buoyant jet}) di dalam
sebuah tangki: air hangat yang masuk dan naik karena lebih ringan, lalu bercampur dengan air di
sekitarnya. Kasus ini menyentuh dua kegagalan di atas sekaligus. Pancaran yang naik membentuk lapisan
temperatur yang tajam, dan setiap difusi tambahan dari pemindahan antar sel mengaburkan lapisan itu.
Ketika dinding tangki kehilangan panas ke lingkungan, alirannya juga tidak lagi berulang sempurna,
karena temperatur rata-rata terus turun.

Benang merah pekerjaannya adalah mengambil lebih banyak informasi dari simulasi penuh dan menyimpannya
di dalam basis data rCFD, alih-alih memakai satu konstanta global. Salah satu contohnya adalah
difusivitas lokal yang direkam dari CFD dan dipakai kembali saat replay, gagasan yang sejalan dengan
koreksi \citet{du2020}. Contoh lain adalah difusivitas negatif yang terkendali untuk mempertajam kembali
gradien yang terlalu kabur, dengan pembatas seperti pada skema flux-corrected transport agar tidak
muncul osilasi. Hasil-hasil kuantitatifnya ada di laporan tugas akhir. Yang ingin saya sampaikan di sini
lebih sederhana: setiap perbaikan mengikuti pola yang sama dengan yang diajarkan di kuliah, yaitu tebakan
yang berani diikuti koreksi berbasis fisika.

\section{rCFD dari sudut pandang machine learning}

Setelah membaca bab-bab sebelumnya, rCFD bisa dilihat dengan mata machine learning. Basis data bingkai
adalah ingatan, recurrence matrix adalah ukuran kemiripan antar keadaan, dan penjahitan adalah mencari
tetangga terdekat di ruang keadaan lalu melompat ke sana. Tidak ada jaringan saraf dan tidak ada
pelatihan dengan gradien. rCFD adalah model pengganti nonparametrik, mirip $k$-nearest neighbour
(\cref{ch:unsup}) yang diterapkan pada lintasan dinamika. Matriks pemindahan antar sel adalah operator
transfer yang dipelajari dari data, saudara dekat DMD dan Koopman (\cref{ch:penemuan}). Dan koreksi neraca
massa adalah contoh paling jelas dari prinsip model hibrida (\cref{ch:hibrida}): biarkan data memberi
dinamika, dan biarkan hukum kekekalan menjaga agar hasilnya tetap fisis.

Kekuatannya juga kelemahannya. Karena hanya memutar ulang dan menyusun kembali apa yang pernah dilihat,
rCFD tidak bisa menghasilkan keadaan aliran yang benar-benar baru. Untuk aliran yang memang mengulang
dirinya, itu justru jaminan: hasilnya tidak akan berhalusinasi. Untuk aliran yang berubah keadaan
operasinya, dibutuhkan basis data dengan banyak keadaan dan cara memilih di antaranya, persis pertanyaan
yang diangkat kuliah tentang proses nyata yang punya lebih dari satu keadaan operasi. Di sinilah rCFD dan model pengganti berbasis jaringan saraf,
seperti yang disinggung di slide terakhir kuliah, mungkin akan saling melengkapi.

\section{Perkembangan riset}

Riset rCFD terus berjalan setelah kuliah. Beberapa publikasi terbaru memperlihatkan arah-arahnya. \citet{hoorijani2024}
mengusulkan rCFD daring, yang menjalankan analisis pengulangan sambil simulasi berjalan, untuk aliran multifase.
\citet{pirker2025seg} menerapkan CFD-DEM dan rCFD pada segregasi ukuran partikel di unggun terfluidisasi gas--padat dengan
dua dan lima ukuran partikel. \citet{lumetzberger2025} mengusulkan aproksimasi momen-propagator untuk transport spesies dalam
aliran turbulen, sebuah cara lain merumuskan langkah transport rCFD. \citet{shariatifar2026} menilai akurasi rCFD untuk
sistem aliran kontinu, dan \citet{shaik2026} membandingkan simulasi CFD-DEM sebuah unggun terfluidisasi bergelembung dengan
pengukuran MRI sebelum memakai rCFD untuk ekstrapolasi waktu. Arah-arah ini sejalan dengan tiga kegagalan di bab ini:
menjaga neraca, memperhalus difusi, dan memperluas rCFD ke aliran yang tidak periodik sempurna.

\sumberbab{Bab ini mengikuti pertemuan tentang recurrence CFD dalam kuliah Numerical Methods in Fluid
Mechanics dan pengalaman selama tugas akhir. Makalah utamanya: \citet{lichtenegger2016} untuk rCFD
berbasis aliran, \citet{pirker2018} untuk rCFD berbasis transport, dan \citet{du2020} untuk penerapan
pada lingkungan bangunan.}
